% GENERATED FILE. Edit canonical lessons, then run npm run generate:books.
% canonical-source-hash: fnv1a64:507c87eb1ca8cc6d
% canonical-lessons: BN-C426-thanda, BN-C426-tushar, BN-C426-jome-jaoya, BN-C426-konkone, BN-C426-gole-jaoya

\chapter{Cold, Frost and Thaw}
\label{ch:bn-cold-frost-and-thaw}
\hlchaptermodality{426}

\begin{chapteropening}
\textbf{By the end of this chapter:} \emph{I can say it is bitterly cold, that water froze and ice melted.}

Complete the last lesson of chapter 426: I can say it is bitterly cold, that water froze and ice melted.
\end{chapteropening}

\section[\texorpdfstring{\bn{ঠান্ডা} (ṭhānḍā)}{ঠান্ডা (ṭhānḍā)}]{\bn{ঠান্ডা} (ṭhānḍā) — cold; chilly; cold weather}

\label{lesson:BN-C426-thanda}



\begin{quote}
Before the new one: say the Bengali for a lightning strike, then the Bengali for to spread over.
\end{quote}

\subsection*{You'll want to know: \bn{ঠান্ডা}}
\textbf{\bn{ঠান্ডা}} — \emph{ṭhānḍā} — \textquotedblleft{}cold; chilly; cold weather\textquotedblright{}.

\begin{quote}
\bn{ঠান্ডা} \bn{জলে} \bn{স্নান} \bn{করে} \bn{অরুণ} \bn{কাঁপছিল।}
\end{quote}

\begin{quote}
\emph{ṭhānḍā jôle snān kore ôrun kā̃pchhilô.}
\end{quote}

\begin{quote}
After bathing in cold water Arun was shivering.
\end{quote}

Also a cold you catch: \bn{ঠান্ডা} \bn{লাগা}. \bn{শীত} is the cold season and the cold you feel; \bn{ঠান্ডা} is cold things too.

\begin{cousinweb}
\textbf{\bn{ঠান্ডা}} (\emph{ṭhānḍā}) is inherited: CDIAL lists Bengali \emph{ṭhāṇḍā}, \textquotedblleft{}cold\textquotedblright{}, under Sanskrit \emph{stabdha}, \textquotedblleft{}stiffened, numb\textquotedblright{} (CDIAL 13676): cold is what makes you stiff.
\end{cousinweb}

\begin{grammarlens}[title={What you've built}]
One more word for getting around.
\end{grammarlens}

\subsection*{Your turn}
\begin{itemize}
\raggedright
  \item \emph{Say it:} \emph{ṭhānḍā}
  \item \emph{Say it:} \emph{ṭhānḍā}, once more
  \item \emph{Say it:} \emph{ṭhānḍā jôle snān kore ôrun kā̃pchhilô.} — After bathing in cold water Arun was shivering.
  \item \emph{Recall:} say the Bengali for a lightning strike, then the Bengali for to spread over, then say \emph{ṭhānḍā} again
\end{itemize}

\begin{tcolorbox}[breakable,colback=teal!4,colframe=teal!35!black,title={Before you move on}]
What does \bn{ঠান্ডা} mean? (\textquotedblleft{}Cold; chilly; cold weather\textquotedblright{}.) Say it once more.
\end{tcolorbox}
\section[\texorpdfstring{\bn{তুষার} (tushār)}{তুষার (tushār)}]{\bn{তুষার} (tushār) — snow; frost}

\label{lesson:BN-C426-tushar}



\begin{quote}
Before the new one: say the Bengali for to spread over, then the Bengali for cold (to the touch).
\end{quote}

\subsection*{You'll want to know: \bn{তুষার}}
\textbf{\bn{তুষার}} — \emph{tushār} — \textquotedblleft{}snow; frost\textquotedblright{}.

\begin{quote}
\bn{শীতে} \bn{উঁচু} \bn{পাহাড়ে} \bn{তুষার} \bn{পড়ে।}
\end{quote}

\begin{quote}
\emph{shīte ũchu pāhāṛe tushār pôṛe.}
\end{quote}

\begin{quote}
In winter snow falls on the high mountains.
\end{quote}

\bn{বরফ} is ice, and snow in everyday talk; \bn{তুষার} is the careful word for snow and frost.

\begin{cousinweb}
\textbf{\bn{তুষার}} (\emph{tushār}) is Sanskrit \emph{tuṣāra}, \textquotedblleft{}cold, frigid\textquotedblright{}, and as a noun \textquotedblleft{}frost, snow\textquotedblright{} (Monier-Williams), taken over whole.
\end{cousinweb}

\begin{grammarlens}[title={What you've built}]
One more word for getting around.
\end{grammarlens}

\subsection*{Your turn}
\begin{itemize}
\raggedright
  \item \emph{Say it:} \emph{tushār}
  \item \emph{Say it:} \emph{tushār}, once more
  \item \emph{Say it:} \emph{shīte ũchu pāhāṛe tushār pôṛe.} — In winter snow falls on the high mountains.
  \item \emph{Recall:} say the Bengali for to spread over, then the Bengali for cold (to the touch), then say \emph{tushār} again
\end{itemize}

\begin{tcolorbox}[breakable,colback=teal!4,colframe=teal!35!black,title={Before you move on}]
What does \bn{তুষার} mean? (\textquotedblleft{}Snow; frost\textquotedblright{}.) Say it once more.
\end{tcolorbox}
\section[\texorpdfstring{\bn{জমে} \bn{যাওয়া} (jôme jāoyā)}{জমে যাওয়া (jôme jāoyā)}]{\bn{জমে} \bn{যাওয়া} (jôme jāoyā) — to freeze; to set hard}

\label{lesson:BN-C426-jome-jaoya}



\begin{quote}
Before the new one: say the Bengali for cold (to the touch), then the Bengali for snow.
\end{quote}

\subsection*{You'll want to know: \bn{জমে} \bn{যাওয়া}}
\textbf{\bn{জমে} \bn{যাওয়া}} — \emph{jôme jāoyā} — \textquotedblleft{}to freeze; to set hard\textquotedblright{}.

\begin{quote}
\bn{পাহাড়ে} \bn{রাতে} \bn{কলের} \bn{জল} \bn{জমে} \bn{গেল।}
\end{quote}

\begin{quote}
\emph{pāhāṛe rāte kôler jôl jôme gelô.}
\end{quote}

\begin{quote}
In the hills the water in the tap froze overnight.
\end{quote}

Past (he/she …): \bn{সেই} \bn{জল} \bn{জমে} \bn{গেল}, \emph{sei jôl jôme gelô}.

Also of a crowd or a party: \bn{আড্ডা} \bn{জমে} \bn{গেল}, the chat got going.

\begin{cousinweb}
\textbf{\bn{জমে} \bn{যাওয়া}} (\emph{jôme jāoyā}) joins the -\bn{এ} form of \bn{জমা}, which CDIAL lists as Bengali \emph{jamā}, \textquotedblleft{}to collect, to coagulate, to freeze\textquotedblright{}, under Sanskrit \emph{yamyátē} (CDIAL 10428), to \bn{যাওয়া}, to go: to go hard.
\end{cousinweb}

\begin{grammarlens}[title={What you've built}]
One more word for getting around.
\end{grammarlens}

\subsection*{Your turn}
\begin{itemize}
\raggedright
  \item \emph{Say it:} \emph{jôme jāoyā}
  \item \emph{Say it:} \emph{jôme jāoyā}, once more
  \item \emph{Say it:} \emph{pāhāṛe rāte kôler jôl jôme gelô.} — In the hills the water in the tap froze overnight.
  \item \emph{Recall:} say the Bengali for cold (to the touch), then the Bengali for snow, then say \emph{jôme jāoyā} again
\end{itemize}

\begin{tcolorbox}[breakable,colback=teal!4,colframe=teal!35!black,title={Before you move on}]
What does \bn{জমে} \bn{যাওয়া} mean? (\textquotedblleft{}To freeze; to set hard\textquotedblright{}.) Say it once more.
\end{tcolorbox}
\section[\texorpdfstring{\bn{কনকনে} (kônkône)}{কনকনে (kônkône)}]{\bn{কনকনে} (kônkône) — bitterly cold; biting (of cold)}

\label{lesson:BN-C426-konkone}



\begin{quote}
Before the new one: say the Bengali for snow, then the Bengali for to freeze.
\end{quote}

\subsection*{You'll want to know: \bn{কনকনে}}
\textbf{\bn{কনকনে}} — \emph{kônkône} — \textquotedblleft{}bitterly cold; biting (of cold)\textquotedblright{}.

\begin{quote}
\bn{কনকনে} \bn{ঠান্ডায়} \bn{আমার} \bn{হাত} \bn{জমে} \bn{গেল।}
\end{quote}

\begin{quote}
\emph{kônkône ṭhānḍāy āmār hāt jôme gelô.}
\end{quote}

\begin{quote}
In the biting cold my hands went numb.
\end{quote}

Only of cold: \bn{কনকনে} \bn{ঠান্ডা}, \bn{কনকনে} \bn{হাওয়া}.

\begin{cousinweb}
\textbf{\bn{কনকনে}} (\emph{kônkône}) is an expressive word whose sound carries the sting of the cold; its origin is not settled.
\end{cousinweb}

\begin{grammarlens}[title={What you've built}]
One more word for getting around.
\end{grammarlens}

\subsection*{Your turn}
\begin{itemize}
\raggedright
  \item \emph{Say it:} \emph{kônkône}
  \item \emph{Say it:} \emph{kônkône}, once more
  \item \emph{Say it:} \emph{kônkône ṭhānḍāy āmār hāt jôme gelô.} — In the biting cold my hands went numb.
  \item \emph{Recall:} say the Bengali for snow, then the Bengali for to freeze, then say \emph{kônkône} again
\end{itemize}

\begin{tcolorbox}[breakable,colback=teal!4,colframe=teal!35!black,title={Before you move on}]
What does \bn{কনকনে} mean? (\textquotedblleft{}Bitterly cold; biting (of cold)\textquotedblright{}.) Say it once more.
\end{tcolorbox}
\section[\texorpdfstring{\bn{গলে} \bn{যাওয়া} (gôle jāoyā)}{গলে যাওয়া (gôle jāoyā)}]{\bn{গলে} \bn{যাওয়া} (gôle jāoyā) — to melt}

\label{lesson:BN-C426-gole-jaoya}



\begin{quote}
Before the new one: say the Bengali for to freeze, then the Bengali for bitterly cold.
\end{quote}

\subsection*{You'll want to know: \bn{গলে} \bn{যাওয়া}}
\textbf{\bn{গলে} \bn{যাওয়া}} — \emph{gôle jāoyā} — \textquotedblleft{}to melt\textquotedblright{}.

\begin{quote}
\bn{রোদে} \bn{বরফ} \bn{গলে} \bn{গেল।}
\end{quote}

\begin{quote}
\emph{rode bôrôph gôle gelô.}
\end{quote}

\begin{quote}
The ice melted in the sun.
\end{quote}

Past (he/she …): \bn{সেই} \bn{বরফ} \bn{গলে} \bn{গেল}, \emph{sei bôrôph gôle gelô}.

Not \bn{গলা}, the throat: the same letters, a different word.

\begin{cousinweb}
\textbf{\bn{গলে} \bn{যাওয়া}} (\emph{gôle jāoyā}) joins the -\bn{এ} form of \bn{গলা}, which CDIAL lists as Bengali \emph{galā}, \textquotedblleft{}to melt, to escape through a hole\textquotedblright{}, under Sanskrit \emph{galati}, \textquotedblleft{}drips, melts\textquotedblright{} (CDIAL 4074), to \bn{যাওয়া}, to go.
\end{cousinweb}

\begin{grammarlens}[title={What you've built}]
One more word for getting around.
\end{grammarlens}

\subsection*{Your turn}
\begin{itemize}
\raggedright
  \item \emph{Say it:} \emph{gôle jāoyā}
  \item \emph{Say it:} \emph{gôle jāoyā}, once more
  \item \emph{Say it:} \emph{rode bôrôph gôle gelô.} — The ice melted in the sun.
  \item \emph{Recall:} say the Bengali for to freeze, then the Bengali for bitterly cold, then say \emph{gôle jāoyā} again
\end{itemize}

\begin{tcolorbox}[breakable,colback=teal!4,colframe=teal!35!black,title={Before you move on}]
What does \bn{গলে} \bn{যাওয়া} mean? (\textquotedblleft{}To melt\textquotedblright{}.) Say it once more.
\end{tcolorbox}
